% preamble-exam-zh.tex — 极简讲解页共享 preamble
% 目标：复刻「题目独立、提示轻框、路线箭头、主体双栏、答案轻收束」的讲义风格。

\documentclass{exam-zh}

% ---------- 基础配置 ----------
\examsetup{
  page/size = a4paper,
  page/show-head = false,
  page/show-foot = false,
  sealline/show = false,
  font = none,
  math-font = none,
}

% ---------- 包 ----------
\usepackage{geometry}
\geometry{top=10mm,bottom=14mm,left=12mm,right=10mm}
\AtBeginDocument{\newgeometry{top=10mm,bottom=14mm,left=12mm,right=10mm}}

\usepackage{xcolor}
\usepackage{needspace}
\usepackage{enumitem}
\usepackage{array}
\usepackage{tabularx}
\usepackage{tcolorbox}
\tcbuselibrary{skins,breakable}
\usepackage{paracol}
\usepackage{graphicx}

% ---------- 打印/屏幕颜色切换 ----------
% 用法：\documentclass[monochrome]{exam-zh} 可降级为灰度。
% 注意：不要使用 \if@xxx 放在 makeatother 外部；这里改成普通条件名。
\newif\ifeduprintcolor
\eduprintcolortrue
\makeatletter
\@ifclasswith{exam-zh}{monochrome}{\eduprintcolorfalse}{}
\makeatother

% ---------- 语义颜色 ----------
\ifeduprintcolor
  \definecolor{edu-blue}{RGB}{31,62,126}
  \definecolor{edu-blue-soft}{RGB}{232,238,250}
  \definecolor{edu-blue-bg}{RGB}{248,250,254}
  \definecolor{edu-ink}{RGB}{30,35,45}
  \definecolor{edu-muted}{RGB}{118,124,136}
  \definecolor{edu-line}{RGB}{209,218,235}
  \definecolor{edu-soft-bg}{RGB}{250,251,253}
  \definecolor{edu-red}{RGB}{168,58,58}
  \definecolor{edu-red-bg}{RGB}{255,247,245}
\else
  \definecolor{edu-blue}{gray}{0.15}
  \definecolor{edu-blue-soft}{gray}{0.90}
  \definecolor{edu-blue-bg}{gray}{0.98}
  \definecolor{edu-ink}{gray}{0.10}
  \definecolor{edu-muted}{gray}{0.45}
  \definecolor{edu-line}{gray}{0.82}
  \definecolor{edu-soft-bg}{gray}{0.97}
  \definecolor{edu-red}{gray}{0.28}
  \definecolor{edu-red-bg}{gray}{0.97}
\fi
\colorlet{edu-diagram-muted}{edu-muted}

% ---------- 全局排版 ----------
\setlength{\parindent}{0pt}
\setlength{\parskip}{0pt}
\linespread{1.16}
\renewcommand{\arraystretch}{1.15}

% ctex 通常提供 \kaishu；这里用它做教师旁白感。
\newcommand{\edukai}{\kaishu}

% ---------- 顶部元信息 ----------
\newcommand{\eduTopBar}[2]{%
  \noindent
  \begin{tabularx}{\linewidth}{@{}Xr@{}}
    {\color{edu-blue}\bfseries #1} &
    {\color{edu-blue}\bfseries #2}
  \end{tabularx}
  \vspace{8mm}
}

% ---------- 轻标题体系 ----------
% 只有真正的大结构才用它；不要给提示/路线滥用标题。
\newcommand{\eduSectionTitle}[1]{%
  \needspace{5\baselineskip}%
  \vspace{2mm}
  {\color{edu-blue}\bfseries\large #1}\par
  \vspace{3mm}
}

\newcommand{\eduSubTitle}[1]{%
  \needspace{4\baselineskip}%
  \vspace{1mm}
  {\color{edu-blue}\bfseries #1}\par
  \vspace{1mm}
}

% ---------- 小问题干：复用 exam (1)(2)(3) 格式 ----------
\newenvironment{eduSubQuestionItem}[1]{%
  \needspace{4\baselineskip}%
  \vspace{2mm}
  \par\noindent
  \hangindent=8mm\hangafter=1
  {\color{edu-blue}\bfseries #1}\enspace
  \begingroup
  \color{edu-ink}
}{%
  \endgroup
  \par
  \vspace{3mm}
}

% ---------- 辅助信息条（解题流程/关键想法/思考）楷体 ----------
\newenvironment{eduAuxItem}[1]{%
  \needspace{4\baselineskip}%
  \vspace{2mm}
  \par\noindent
  \hangindent=8mm\hangafter=1
  {\color{edu-blue}\edukai $\triangleright$~#1}\enspace
  \begingroup
  \color{edu-ink}
  \edukai
}{%
  \endgroup
  \par
  \vspace{4mm}
}

\newcommand{\eduColumnTitle}[2]{%
  {\color{edu-blue}\bfseries #1\hspace{1.5mm}#2}\par
  \vspace{2.5mm}
}

% ---------- 图标：不用外部 icon 字体，保证可移植 ----------
\newcommand{\eduIconBulb}{\(\circ\)}
\newcommand{\eduIconBook}{\(\square\)}
\newcommand{\eduIconCheck}{\(\checkmark\)}
\newcommand{\eduIconPin}{\(\diamond\)}
\newcommand{\eduIconNote}{\(\triangleright\)}

% ---------- 题目区 ----------
% 题目区是页面入口：弱标签 + 一条长蓝线 + 大题干。
\newenvironment{eduProblemBlock}[1][题目]{%
  \needspace{8\baselineskip}
  {\small\color{edu-muted}#1}\par
  \vspace{1.5mm}
  {\color{edu-blue}\hrule height 0.55pt}
  \vspace{4.5mm}
  \begingroup
  \color{edu-ink}
  \large
  \setlength{\parskip}{2.2mm}
  \linespread{1.20}\selectfont
}{%
  \par
  \endgroup
  \vspace{6mm}
}

\newenvironment{eduSubQuestions}{%
  \begin{enumerate}[
    label=(\arabic*),
    leftmargin=8mm,
    itemsep=2.0mm,
    topsep=2.0mm,
    parsep=0pt
  ]
}{%
  \end{enumerate}
}

% ---------- 读题提示：轻框，不做同级标题 ----------
\newtcolorbox{eduReadTipBox}{
  enhanced,
  breakable,
  colback=white,
  colframe=edu-blue!45,
  boxrule=0.45pt,
  arc=1.6mm,
  left=10mm,
  right=4mm,
  top=5mm,
  bottom=5mm,
  before skip=1mm,
  after skip=7mm,
  fontupper=\edukai\normalsize,
  colupper=edu-ink,
  overlay={
    \node[anchor=north west, text=edu-blue] at ([xshift=3.2mm,yshift=-3.2mm]frame.north west)
    {\small\eduIconNote};
  }
}

\newenvironment{eduTipList}{%
  \begin{itemize}[
    label=\textbullet,
    leftmargin=5mm,
    itemsep=1.2mm,
    topsep=0pt,
    parsep=0pt
  ]
}{%
  \end{itemize}
}

% ---------- 解题路线：虚线框 + 文字箭头 ----------
\newenvironment{eduRouteFlow}{%
  \needspace{4\baselineskip}%
  \vspace{1mm}
  \begin{center}
  \normalsize
}{%
  \end{center}
  \vspace{7mm}
}
\newcommand{\eduRouteStep}[1]{%
  {\color{edu-ink}#1}%
}
\newcommand{\eduRouteArrow}{%
  \;$\boldsymbol{\to}$\;%
}

% ---------- 主体讲解：使用 paracol 实现跨页自适应双栏 ----------
\newenvironment{eduDualExplain}{%
  \needspace{6\baselineskip}% 防止标题孤行
  \columnratio{0.26, 0.74}%
  \setlength{\columnsep}{8mm}%
  \columnseprule=0.45pt%
  \colseprulecolor{edu-blue}%
  \begin{paracol}{2}% 开启双栏
}{%
  \end{paracol}% 结束双栏
  \vspace{5mm}
}

\newenvironment{eduExplainLeft}{%
  \setlength{\linewidth}{\columnwidth}%
  \hsize=\columnwidth
  \normalsize\color{edu-ink}%
}{%
  \switchcolumn
}

\newcommand{\eduColumnDivider}{}

\newenvironment{eduExplainRight}{%
  \setlength{\linewidth}{\columnwidth}%
  \hsize=\columnwidth
  \normalsize\color{edu-ink}%
}{%
}

\newtcolorbox{eduSideHintBox}[1]{
  enhanced,
  breakable,
  colback=edu-blue-bg,
  colframe=edu-line,
  boxrule=0.35pt,
  borderline west={1.5pt}{0pt}{edu-blue},
  arc=0.8mm,
  left=2.4mm,
  right=2.4mm,
  top=2mm,
  bottom=2mm,
  before skip=1mm,
  after skip=1.5mm,
  title={\color{edu-blue}\bfseries\small #1},
  coltitle=edu-blue,
  attach title to upper={\par\vspace{0.7mm}},
  fontupper=\small
}

\newtcolorbox{eduSideMistakeBox}[1]{
  enhanced,
  breakable,
  colback=edu-red-bg,
  colframe=edu-red!35,
  boxrule=0.35pt,
  borderline west={1.5pt}{0pt}{edu-red},
  arc=0.8mm,
  left=2.4mm,
  right=2.4mm,
  top=2mm,
  bottom=2mm,
  before skip=1mm,
  after skip=1.5mm,
  title={\color{edu-red}\bfseries\small #1},
  coltitle=edu-red,
  attach title to upper={\par\vspace{0.7mm}},
  fontupper=\small
}

\newtcolorbox{eduSideNoteBox}[1]{
  enhanced,
  breakable,
  colback=edu-soft-bg,
  colframe=edu-line,
  boxrule=0.35pt,
  borderline west={1.5pt}{0pt}{edu-muted},
  arc=0.8mm,
  left=2.4mm,
  right=2.4mm,
  top=2mm,
  bottom=2mm,
  before skip=1mm,
  after skip=1.5mm,
  title={\color{edu-muted}\bfseries\small #1},
  coltitle=edu-muted,
  attach title to upper={\par\vspace{0.7mm}},
  fontupper=\small
}

\newcommand{\eduSideHint}[2]{%
  \begin{eduSideHintBox}{#1}#2\end{eduSideHintBox}%
}
\newcommand{\eduSideMistake}[2]{%
  \begin{eduSideMistakeBox}{#1}#2\end{eduSideMistakeBox}%
}
\newcommand{\eduSideNote}[2]{%
  \begin{eduSideNoteBox}{#1}#2\end{eduSideNoteBox}%
}

\newcommand{\eduExplainStep}[3]{%
  \needspace{3\baselineskip}%
  \vspace{1.2mm}%
  {\color{edu-blue}\bfseries step#1.\enspace #2}\par
  \vspace{0.8mm}%
  \begingroup
  \color{edu-ink}#3\par
  \endgroup
  \vspace{2.2mm}%
}

\newenvironment{eduNumberList}{%
  \begin{enumerate}[
    label=\arabic*.,
    leftmargin=6mm,
    itemsep=4.8mm,
    topsep=0pt,
    parsep=0pt
  ]
}{%
  \end{enumerate}
}

\newenvironment{eduBulletList}{%
  \begin{itemize}[
    label=\textbullet,
    leftmargin=5mm,
    itemsep=1.2mm,
    topsep=0pt,
    parsep=0pt
  ]
}{%
  \end{itemize}
}

% ---------- 底部双栏：答案 + 方法提醒 ----------
\newenvironment{eduBottomGrid}{%
  \columnratio{0.32, 0.68}%
  \setlength{\columnsep}{11mm}%
  \columnseprule=0pt%
  \begin{paracol}{2}%
}{%
  \end{paracol}%
}

\newenvironment{eduSummaryLeft}{%
}{%
}

\newcommand{\eduSummaryDivider}{%
  \switchcolumn
}

\newenvironment{eduSummaryRight}{%
}{%
}

% ---------- 答案/方法提醒：轻量收束 ----------
\newtcolorbox{eduSummaryBlock}[2]{
  enhanced,
  breakable,
  colback=edu-blue-bg,
  colframe=edu-line,
  boxrule=0.45pt,
  arc=1.2mm,
  left=3.5mm,
  right=3.5mm,
  top=3mm,
  bottom=3mm,
  before skip=1mm,
  after skip=2.5mm,
  title={\color{edu-blue}\bfseries #1\hspace{1.5mm}#2},
  coltitle=edu-blue,
  attach title to upper={\par\vspace{1.5mm}},
  fontupper=\normalsize
}

% ---------- 例题单栏解答卡片 ----------
\newtcolorbox{eduSolutionBox}[1]{
  enhanced,
  breakable,
  colback=white,
  colframe=edu-blue!40,
  boxrule=0.55pt,
  arc=1.2mm,
  left=4mm,
  right=4mm,
  top=3.5mm,
  bottom=3.5mm,
  before skip=2mm,
  after skip=3mm,
  title={\color{edu-blue}\bfseries \eduIconBook\hspace{1.5mm}#1},
  coltitle=edu-blue,
  attach title to upper={\par\vspace{1.5mm}},
  fontupper=\normalsize
}

% ---------- 变式训练：完整题目 + 学生留白 ----------
\newtcolorbox{eduVariationTraining}[1]{
  enhanced,
  breakable,
  colback=white,
  colframe=edu-blue!45,
  boxrule=0.5pt,
  borderline west={1.8pt}{0pt}{edu-blue},
  arc=1.2mm,
  left=4mm,
  right=4mm,
  top=3.2mm,
  bottom=3.2mm,
  before skip=2mm,
  after skip=4mm,
  title={\color{edu-blue}\bfseries #1},
  coltitle=edu-blue,
  attach title to upper={\par\vspace{1.8mm}},
  fontupper=\normalsize,
  colupper=edu-ink
}

\newenvironment{eduPlainList}{%
  \begin{enumerate}[
    label=(\arabic*),
    leftmargin=7mm,
    itemsep=1.2mm,
    topsep=0pt,
    parsep=0pt
  ]
}{%
  \end{enumerate}
}

% ---------- 兼容旧 block 的轻量盒子 ----------
\newtcolorbox{eduSoftBox}{
  enhanced,
  breakable,
  colback=edu-soft-bg,
  colframe=edu-line,
  boxrule=0.35pt,
  arc=1mm,
  left=3mm,
  right=3mm,
  top=2mm,
  bottom=2mm,
  before skip=1mm,
  after skip=2mm,
  fontupper=\small
}

\newtcolorbox{eduMistakeBox}{
  enhanced,
  breakable,
  colback=edu-red-bg,
  colframe=edu-red!40,
  boxrule=0.35pt,
  arc=1mm,
  left=3mm,
  right=3mm,
  top=2.5mm,
  bottom=2.5mm,
  before skip=1mm,
  after skip=2mm,
  fontupper=\normalsize
}

\newtcolorbox{eduLegacySolution}{
  enhanced,
  breakable,
  colback=edu-soft-bg,
  colframe=edu-line,
  boxrule=0.35pt,
  arc=1mm,
  left=3mm,
  right=3mm,
  top=2mm,
  bottom=2mm,
  before skip=-2mm,
  after skip=4mm,
  fontupper=\small
}

\newtcolorbox{eduTeacherNote}{
  enhanced,
  breakable,
  colback=edu-soft-bg,
  colframe=edu-line,
  boxrule=0pt,
  borderline west={1.5pt}{0pt}{edu-muted},
  left=2.5mm,
  right=2.5mm,
  top=2mm,
  bottom=2mm,
  before skip=1mm,
  after skip=2mm,
  fontupper=\footnotesize,
  colupper=edu-muted
}

% ---------- 答题区 ----------
\newcommand{\answerarea}[1][40mm]{%
  \vspace{3mm}%
  \begin{tcolorbox}[
    enhanced,
    breakable,
    colback=white,
    colframe=edu-line,
    boxrule=0.55pt,
    arc=0.8mm,
    height=#1,
    valign=top,
  ]
  \end{tcolorbox}%
}

\newcommand{\diagramcolfigure}[3]{%
  \begin{minipage}[t]{#2}
    \vspace{0pt}\centering
    \includegraphics[width=\linewidth]{\detokenize{#1}}
    \if\relax\detokenize{#3}\relax\else
      \par\vspace{1mm}{\scriptsize\color{edu-diagram-muted}#3}
    \fi
  \end{minipage}%
}

\newcommand{\diagramrowitem}[4]{%
  \begin{minipage}[t]{#2}
    \vspace{0pt}\centering
    \includegraphics[width=\linewidth]{\detokenize{#1}}
    \if\relax\detokenize{#3}\relax\else
      \par\vspace{1mm}{\scriptsize\bfseries #3}
    \fi
    \if\relax\detokenize{#4}\relax\else
      \par\vspace{0.5mm}{\scriptsize\color{edu-diagram-muted}#4}
    \fi
  \end{minipage}%
}

\newcommand{\answerareawithdiagramcol}[4]{%
  \vspace{3mm}%
  \begin{tcolorbox}[
    enhanced,
    breakable,
    colback=white,
    colframe=edu-line,
    boxrule=0.55pt,
    arc=0.8mm,
    height=#1,
    valign=top,
  ]
  \begin{minipage}[t]{\dimexpr\linewidth-#3-6mm\relax}
    \vspace{0pt}
  \end{minipage}\hfill
  \diagramcolfigure{#2}{#3}{#4}
  \end{tcolorbox}%
}

\newcommand{\answerareawithdiagram}[4]{%
  \answerareawithdiagramcol{#1}{#2}{#3}{#4}%
}

\examsetup{
  question/show-answer = true,
  fillin/show-answer = true,
  solution/show-solution = show-stay,
}

\begin{document}

\clearpage
\eduSectionTitle{模型一：直角顶点已知，求等腰直角第三点}

\begin{eduProblemBlock}[例题 1]
在平面直角坐标系中，等腰直角三角形 $ABC$ 满足 $\angle A=90^\circ$，$AB=AC$。
已知 $A(-2,3)$，$B(4,7)$，点 $C$ 在点 $A$ 的右下方，求点 $C$ 的坐标。

\end{eduProblemBlock}





\begin{eduAuxItem}{思路导航}
读出 $AB$ 的走法$\;\to\;$把走法旋转成 $AC$$\;\to\;$根据位置取方向$\;\to\;$由 $A$ 到 $C$ 定坐标\end{eduAuxItem}




\begin{eduSubQuestionItem}{例题 1}
求底角顶点 $C$ 的坐标。\end{eduSubQuestionItem}

\begin{eduDualExplain}
  \begin{eduExplainLeft}
    \eduColumnTitle{\eduIconBulb}{小贴士}
        \eduSideHint{为什么能交换步数}{过 $A$ 作竖线，过 $B,C$ 向这条竖线作垂线，会得到两个全等的小直角三角形，所以横竖步数交换。}
        \eduSideMistake{别漏另一个方向}{如果题目没有说 $C$ 在哪一侧，一般要写两个候选点。}
        \eduSideHint{先说走法}{不要一上来写公式，先把 $AB$ 说成“右几上几”。}
  \end{eduExplainLeft}

  \begin{eduExplainRight}
    \eduColumnTitle{\eduIconBook}{解答}
      \eduExplainStep{1}{读出 $AB$ 的走法}{从 $A(-2,3)$ 到 $B(4,7)$：横坐标加 $6$，纵坐标加 $4$，所以走法是“右 $6$ 上 $4$”。}
      \eduExplainStep{2}{把走法旋转成 $AC$}{因为 $AB=AC$ 且 $AB\perp AC$，所以把“右 $6$ 上 $4$”变成“右 $4$ 下 $6$”或“左 $4$ 上 $6$”。}
      \eduExplainStep{3}{根据位置取方向}{题目说 $C$ 在 $A$ 的右下方，所以取“右 $4$ 下 $6$”。}
      \eduExplainStep{4}{由 $A$ 到 $C$ 定坐标}{从 $A(-2,3)$ 出发，右 $4$ 得 $x=2$，下 $6$ 得 $y=-3$，所以 $C(2,-3)$。}
    \vspace{1mm}
    {\color{edu-blue}\bfseries\small 检查}\par\vspace{1mm}
    \begin{eduBulletList}
      \item $AB$ 是右 $6$ 上 $4$，$AC$ 是右 $4$ 下 $6$，长度相同，方向垂直。    \end{eduBulletList}
  \end{eduExplainRight}
\end{eduDualExplain}




\begin{eduMistakeBox}
\eduSubTitle{易错提醒}坐标题里最容易把“走法”和“坐标”混在一起。$A(-2,3)$ 到 $B(4,7)$ 的走法是右 $6$ 上 $4$，不是点 $(6,4)$；最后求 $C$ 时还要从 $A$ 出发走一次。
\end{eduMistakeBox}




\begin{eduSummaryBlock}{\eduIconPin}{方法提醒}
\begin{eduBulletList}
  \item 先读 $AB$：右/左几，上/下几。  \item 再旋转：水平竖直交换，并根据图形位置决定方向。  \item 最后从直角顶点出发，按新走法确定第三点。\end{eduBulletList}
\end{eduSummaryBlock}



\clearpage
\eduSectionTitle{模型二：斜边两端已知，求直角顶点}

\begin{eduProblemBlock}[例题 2]
在平面直角坐标系中，等腰直角三角形 $ABC$ 满足 $\angle A=90^\circ$。
已知 $B(-3,4)$，$C(5,-2)$，点 $A$ 在第一象限，求点 $A$ 的坐标。

\end{eduProblemBlock}





\begin{eduAuxItem}{思路导航}
取 $BC$ 中点 $D$$\;\to\;$读出 $DB$ 的走法$\;\to\;$把 $DB$ 旋转成 $DA$$\;\to\;$根据位置确定 $A$\end{eduAuxItem}




\begin{eduSubQuestionItem}{例题 2}
求顶角顶点 $A$ 的坐标。\end{eduSubQuestionItem}

\begin{eduDualExplain}
  \begin{eduExplainLeft}
    \eduColumnTitle{\eduIconBulb}{小贴士}
        \eduSideHint{为什么取中点}{已知的是斜边两端 $B,C$，中点 $D$ 能把未知的直角顶点转化成上一节的“底角顶点”。}
        \eduSideMistake{不要从 $B$ 直接转}{这类题的旋转中心是中点 $D$，不是 $B$ 或 $C$。}
  \end{eduExplainLeft}

  \begin{eduExplainRight}
    \eduColumnTitle{\eduIconBook}{解答}
      \eduExplainStep{1}{取 $BC$ 中点 $D$}{$D$ 的坐标为 $\left(\dfrac{-3+5}{2},\dfrac{4+(-2)}{2}\right)=(1,1)$。}
      \eduExplainStep{2}{读出 $DB$ 的走法}{从 $D(1,1)$ 到 $B(-3,4)$ 是“左 $4$ 上 $3$”。}
      \eduExplainStep{3}{把 $DB$ 旋转成 $DA$}{与“左 $4$ 上 $3$”等长垂直的走法可以是“右 $3$ 上 $4$”，也可以是“左 $3$ 下 $4$”。}
      \eduExplainStep{4}{根据位置确定 $A$}{从 $D(1,1)$ 走“右 $3$ 上 $4$”得 $A(4,5)$；另一个候选为 $A(-2,-3)$。题目说 $A$ 在第一象限，所以 $A(4,5)$。}
    \vspace{1mm}
    {\color{edu-blue}\bfseries\small 检查}\par\vspace{1mm}
    \begin{eduBulletList}
      \item $D$ 是 $BC$ 中点，$DA$ 与 $DB$ 都是右三角形的一半斜边，方向也垂直。    \end{eduBulletList}
  \end{eduExplainRight}
\end{eduDualExplain}




\begin{eduMistakeBox}
\eduSubTitle{易错提醒}已知 $B,C$ 求 $A$ 时，先找中点 $D$。如果跳过 $D$，很容易把整条 $BC$ 当成一条腰，长度会错一倍。
\end{eduMistakeBox}




\begin{eduSummaryBlock}{\eduIconPin}{方法提醒}
\begin{eduBulletList}
  \item 已知两个底角顶点，先取中点。  \item 从中点到一个端点读走法，再旋转 $90^\circ$。  \item 两个候选点用题目给的位置条件筛选。\end{eduBulletList}
\end{eduSummaryBlock}



\clearpage
\eduSectionTitle{模型三：45°方向直线与函数交点}

\begin{eduProblemBlock}[例题 3]
已知 $A(2,5)$，$B(5,6)$。直线 $BC$ 与射线 $BA$ 成 $45^\circ$，且对应的辅助点 $P$ 在点 $A$ 的右下方。
\begin{enumerate}[label=(\arabic*)]
  \item 求直线 $BC$ 的解析式；
  \item 若直线 $BC$ 与反比例函数 $y=\dfrac{30}{x}$ 的图像除 $B$ 外还交于点 $C$，求点 $C$ 的坐标与 $\triangle ABC$ 的面积。
\end{enumerate}

\end{eduProblemBlock}





\begin{eduAuxItem}{思路导航}
读出 $AB$ 的走法$\;\to\;$作等腰直角三角形 $ABP$$\;\to\;$用 $B,P$ 求直线$\;\to\;$联立反比例函数求 $C$$\;\to\;$用坐标底高求面积\end{eduAuxItem}




\begin{eduSubQuestionItem}{(1)}
求直线 $BC$ 的解析式。\end{eduSubQuestionItem}

\begin{eduDualExplain}
  \begin{eduExplainLeft}
    \eduColumnTitle{\eduIconBulb}{小贴士}
        \eduSideHint{为什么找 $P$}{题目只给了 $45^\circ$，直接求直线不好做；补出等腰直角三角形后，$BP$ 就是一条明确的 45 度方向线。}
        \eduSideMistake{方向要看题}{本题说 $P$ 在 $A$ 的右下方，所以只取“右 $1$ 下 $3$”。}
  \end{eduExplainLeft}

  \begin{eduExplainRight}
    \eduColumnTitle{\eduIconBook}{解答}
      \eduExplainStep{1}{读出 $AB$ 的走法}{从 $A(2,5)$ 到 $B(5,6)$ 是“右 $3$ 上 $1$”。}
      \eduExplainStep{2}{作等腰直角三角形 $ABP$}{以 $A$ 为直角顶点，取 $P$ 在 $A$ 的右下方，则 $AP$ 的走法是“右 $1$ 下 $3$”，所以 $P(3,2)$。}
      \eduExplainStep{3}{用 $B,P$ 求直线}{直线 $BP$ 的斜率为 $k=\dfrac{6-2}{5-3}=2$，代入 $B(5,6)$，得 $6=2\cdot5+b$，所以 $b=-4$，即 $BC:y=2x-4$。}
  \end{eduExplainRight}
\end{eduDualExplain}




\begin{eduSubQuestionItem}{(2)}
求点 $C$ 的坐标与 $\triangle ABC$ 的面积。\end{eduSubQuestionItem}

\begin{eduDualExplain}
  \begin{eduExplainLeft}
    \eduColumnTitle{\eduIconBulb}{小贴士}
        \eduSideHint{后续题怎么接}{一旦直线 $BC$ 求出来，后面的交点、面积就是一次函数与反比例函数的综合。}
        \eduSideMistake{面积不要凭图猜}{点不在水平或竖直底边上时，用坐标面积更稳。}
  \end{eduExplainLeft}

  \begin{eduExplainRight}
    \eduColumnTitle{\eduIconBook}{解答}
      \eduExplainStep{1}{联立反比例函数求 $C$}{联立 $y=2x-4$ 与 $y=\dfrac{30}{x}$，得 $2x-4=\dfrac{30}{x}$，即 $x^2-2x-15=0$。解得 $x=5$ 或 $x=-3$。$x=5$ 对应点 $B$，所以另一个交点为 $C(-3,-10)$。}
      \eduExplainStep{2}{用坐标底高求面积}{用坐标面积：$S_{\triangle ABC}=\dfrac12\left|(5-2)(-10-5)-(6-5)(-3-2)\right|=20$。}
  \end{eduExplainRight}
\end{eduDualExplain}




\begin{eduMistakeBox}
\eduSubTitle{易错提醒}$\angle ABC=45^\circ$ 的顶点是 $B$，所以最后要找的是过 $B$ 的直线。辅助三角形中 $A$ 是直角顶点，$P$ 是为了帮我们找方向，不是题目原本给的点。
\end{eduMistakeBox}




\begin{eduSummaryBlock}{\eduIconPin}{方法提醒}
\begin{eduBulletList}
  \item 看到 $\angle ABC=45^\circ$，先让 $A$ 做直角顶点，补等腰直角三角形 $ABP$。  \item 求出 $P$ 后，直线 $BP$ 就是 $BC$ 的候选直线。  \item 若题目没有给方向，要写两条候选直线或说明需要附加条件。\end{eduBulletList}
\end{eduSummaryBlock}



\clearpage
\eduSectionTitle{条件辨析：∠ABC 与 ∠ACB 的差异}

\begin{eduProblemBlock}[例题 4]
已知 $A(1,2)$，$B(7,2)$。只知道 $\angle ACB=45^\circ$，能不能唯一求出直线 $BC$ 的解析式？说明理由。

\end{eduProblemBlock}





\begin{eduAuxItem}{思路导航}
先看角的顶点$\;\to\;$判断是否唯一$\;\to\;$补充条件后才能求\end{eduAuxItem}




\begin{eduSubQuestionItem}{例题 4}
判断是否能唯一求出 $BC$ 的解析式。\end{eduSubQuestionItem}

\begin{eduDualExplain}
  \begin{eduExplainLeft}
    \eduColumnTitle{\eduIconBulb}{小贴士}
        \eduSideHint{先问顶点}{角写在中间的字母就是顶点。$ABC$ 的顶点是 $B$，$ACB$ 的顶点是 $C$。}
        \eduSideMistake{不要套错模型}{上一节的补点 $P$ 是为了解决 $\angle ABC=45^\circ$，不能直接搬到 $\angle ACB=45^\circ$。}
  \end{eduExplainLeft}

  \begin{eduExplainRight}
    \eduColumnTitle{\eduIconBook}{解答}
      \eduExplainStep{1}{先看角的顶点}{$\angle ACB$ 的顶点是 $C$，但 $C$ 的位置没有给出。}
      \eduExplainStep{2}{判断是否唯一}{不同位置的点 $C$ 都可能看见同一条 $AB$ 成 $45^\circ$，所以直线 $BC$ 不唯一。}
      \eduExplainStep{3}{补充条件后才能求}{若再给 $C$ 在某条直线上，或给 $C$ 的横坐标、纵坐标，才能把 $C$ 定下来，再求 $BC$。}
    \vspace{1mm}
    {\color{edu-blue}\bfseries\small 结论}\par\vspace{1mm}
    \begin{eduBulletList}
      \item 不能唯一求出；题目还缺少确定 $C$ 的条件。    \end{eduBulletList}
  \end{eduExplainRight}
\end{eduDualExplain}




\begin{eduMistakeBox}
\eduSubTitle{易错提醒}遇到 $45^\circ$ 题，不要只看数字 $45^\circ$。先看角的顶点是谁，再决定这条 45 度方向线应该过哪个点。
\end{eduMistakeBox}




\begin{eduSummaryBlock}{\eduIconPin}{方法提醒}
\begin{eduBulletList}
  \item 角名中间的字母是顶点。  \item 顶点已知，常能求一条候选直线；顶点未知，先判断条件是否足够。  \item 条件不足时，规范写法是“不能唯一确定，还需要补充条件”。\end{eduBulletList}
\end{eduSummaryBlock}




\end{document}
